
\section*{Ethical Considerations}
We discuss ethical considerations by centering the impacts on each stakeholder across two phases: the \emph{research process} (attack design and evaluation) and the \emph{publication of results} (deployment implications).
% 
We then describe the mitigation and conclude with justification for conducting this research.

\paragraph{Stakeholder Analysis.}
Consider three stakeholder groups:
\begin{icompact}
\item \emph{Model providers and infrastructure operators} who develop LLM serving systems and cache reuse frameworks. Our findings reveal an underexplored threat in position-independent KV-cache reuse, showing that efficiency optimizations can introduce new attack surfaces. This can inform safer cache reuse mechanisms and improve system-level security posture. The potential negative impact is that attack insights could be misused against systems that deploy cache reuse without sufficient integrity protections.

\item \emph{End users} of LLM systems, whose interactions depend on reliable context integrity and isolation. They may benefit from defenses against unintended cross-context influence, while insecure reuse mechanisms could affect response reliability in shared infrastructure settings.

\item \emph{Researchers and practitioners} studying LLM systems, who benefit from understanding system-level risks. Our work frames KV-cache reuse as a security-relevant component and may motivate further defensive research, while detailed technical analysis could lower the barrier for reproducing attacks without appropriate safeguards.
\end{icompact}

\paragraph{Mitigation (Implemented).}
We took steps to minimize harm to the identified stakeholders and promote defensive progress.
\begin{icompact}
\item \emph{For model providers and infrastructure operators, responsible disclosure}: Prior to publication, we contacted the founders of a startup commercializing position-independent cache reuse technology and shared our threat model, technical analysis, and empirical results. Follow-up discussions were scheduled to explore safer reuse designs.

\item \emph{For end users, no real-world harm}: All experiments used public datasets and public or locally deployed models in controlled offline environments. We did not use private user data, interact with real users, or test production services.

\item \emph{For researchers and practitioners, defense analysis}: We evaluate mitigation directions, including recomputation-based protection, cache compression, and existing defensive mechanisms, and provide system-level recommendations rather than only exploit guidance.
\end{icompact}

\paragraph{Recommended Future Safeguards.}
We suggest that real-world systems incorporate safeguards that address the risks faced by both infrastructure operators and end users: (1) a hybrid mechanism that recomputes high-impact KV states while compressing irrelevant entries, ensuring that user-computed data constitutes the majority of the cache; and (2) a deviation-based rejection policy that denies reuse whenever recomputation reveals deviation exceeding a predefined safety threshold.

\paragraph{Justification for Research.}
KV cache reuse improves LLM efficiency but introduces underexplored security risks. Identifying and responsibly disclosing these risks is essential to ensure efficiency gains do not compromise system integrity.

\section*{Open Science}
To support open science, \sys is publicly at \url{https://github.com/YichiCS/KV-Cache-Hijack} and \url{https://zenodo.org/records/20403786}. 

\section*{Acknowledgment} 
We thank the reviewers and shepherd for their constructive comments on our work. 
% 
The Authors acknowledge the National Artificial Intelligence Research Resource (NAIRR) Pilot for contributing to this research result.
% 
Huan Zhang is supported in part by the AI2050 program at Schmidt Sciences (AI2050 Early Career Fellowship).

\bibliographystyle{plain}
\bibliography{reference}